\section{One normalized functional equation for all derivative bridges}
\label{sec:bridges}

The Stieltjes drafts connect derivatives of \(L(s,\chi)\) at a positive
integer with derivatives at its reflected negative integer. A single
identity of analytic germs organizes all these formulas and treats
trivial zeros without differentiating singular logarithms. It also
makes the conjugation required for complex characters unavoidable.
The input is the classical primitive Dirichlet functional equation
and the reflection and duplication formulas for \(\Gamma\)
\cite{DLMF}; the all-order normalized form below is a systematic
continuation of the bridges used in the corpus.

\subsection{Regularization and the normalized germ}

Let \(\chi\) be primitive of conductor \(q\), allowing the trivial
character of conductor \(1\), for which \(L(s,\chi)=\zeta(s)\).
Fix \(k\geq1\), and write
\[
 \delta=\frac{1-\chi(-1)}2\in\{0,1\},\qquad
 \nu=
 \begin{cases}
 1,&\chi(-1)=(-1)^k,\\
 0,&\chi(-1)\ne(-1)^k.
 \end{cases}
\]
Define
\begin{equation}
 B(t)=\frac{L(-k+t,\chi)}{t^\nu},
 \qquad
 B^\#(z)=\overline{B(\overline z)},
 \label{eq:bridge-regularization}
\end{equation}
with the removable value filled in when \(\nu=1\).
The function \(B^\#\) is holomorphic near zero; it is coefficientwise
conjugation of the Taylor series of \(B\).

\begin{theorem}[Normalized positive--negative germ]
\label{thm:bridge-normalized}
The quantities \(L(k+1,\chi)\) and \(B(0)\) are nonzero. On a
sufficiently small disk about \(z=0\),
\begin{equation}
 \frac{L(k+1+z,\chi)}{L(k+1,\chi)}
 =
 K_\nu(z)\frac{B^\#(-z)}{B^\#(0)},
 \label{eq:bridge-normalized-germ}
\end{equation}
where
\begin{equation}
 \begin{aligned}
 K_\nu(z)&=
 \left(\frac{2\pi}{q}\right)^z
 \frac{\Gamma(k+1)}{\Gamma(k+1+z)}\,T_\nu(z),\\
 T_0(z)&=\sec\frac{\pi z}{2},\qquad
 T_1(z)=\frac{\pi z/2}{\sin(\pi z/2)},\quad T_1(0)=1.
 \end{aligned}
 \label{eq:bridge-elementary-kernel}
\end{equation}
All three normalized factors in
\eqref{eq:bridge-normalized-germ} equal \(1\) at the origin.
\end{theorem}

\begin{proof}
The Euler product gives \(L(k+1,\chi)\ne0\), since \(k+1>1\).
Use the completed function
\[
 \Lambda(s,\chi)=
 \left(\frac q\pi\right)^{(s+\delta)/2}
 \Gamma\!\left(\frac{s+\delta}{2}\right)L(s,\chi)
\]
and its functional equation
\begin{equation}
 \Lambda(s,\chi)=\varepsilon_\chi\Lambda(1-s,\overline\chi),
 \qquad |\varepsilon_\chi|=1.
 \label{eq:bridge-completed-fe}
\end{equation}
At \(s=-k\), the gamma factor is finite and nonzero if
\(\nu=0\), and has a simple pole if \(\nu=1\).
The reflected completed value at \(k+1\) is finite and nonzero.
Consequently \(L(-k,\chi)\ne0\) in the first case, and
\(L(s,\chi)\) has a simple zero at \(-k\) in the second.
In either case \(B\) is holomorphic and \(B(0)\ne0\).

Put
\[
 a=\frac{\delta-k}{2},\qquad
 b=\frac{k+1+\delta}{2}.
\]
Solving \eqref{eq:bridge-completed-fe} at \(s=k+1+z\) gives
\begin{equation}
 L(k+1+z,\chi)
 =\varepsilon_\chi
 \left(\frac q\pi\right)^{-k-\frac12-z}
 \frac{\Gamma(a-z/2)}{\Gamma(b+z/2)}
 L(-k-z,\overline\chi).
 \label{eq:bridge-fe-before-normalization}
\end{equation}
Since
\(L(-k-z,\overline\chi)=(-z)^\nu B^\#(-z)\),
the numerator
\[
 F_\nu(z)=(-z)^\nu\Gamma(a-z/2)
\]
has a nonzero removable value at the origin.
Dividing \eqref{eq:bridge-fe-before-normalization} by its value
at \(z=0\) cancels the root number and all constant factors:
\begin{equation}
 \frac{L(k+1+z,\chi)}{L(k+1,\chi)}
 =
 \left(\frac q\pi\right)^{-z}
 \frac{F_\nu(z)}{F_\nu(0)}
 \frac{\Gamma(b)}{\Gamma(b+z/2)}
 \frac{B^\#(-z)}{B^\#(0)}.
 \label{eq:bridge-normalized-intermediate}
\end{equation}

If \(\nu=0\), then \(a\) is a half-integer.
The gamma reflection formula and
\(\sin(\pi(a-z/2))=\sin(\pi a)\cos(\pi z/2)\)
give
\begin{equation}
 \frac{F_0(z)}{F_0(0)}
 =\frac{\Gamma(1-a)}{\Gamma(1-a+z/2)}
       \sec\frac{\pi z}{2}.
 \label{eq:bridge-gamma-half-integer}
\end{equation}
If \(\nu=1\), write \(a=-m\), \(m\geq0\).
Then \(F_1(0)=2(-1)^m/m!\), and the same reflection formula
gives
\begin{equation}
 \frac{F_1(z)}{F_1(0)}
 =\frac{\Gamma(1-a)}{\Gamma(1-a+z/2)}
       \frac{\pi z/2}{\sin(\pi z/2)}.
 \label{eq:bridge-gamma-integer}
\end{equation}
For example, the pole in \(\Gamma(-m-z/2)\) is canceled by
the explicit factor \(-z\), so no logarithm of a function
vanishing at zero has been taken.

The unordered pair \(\{b,1-a\}\) is
\(\{(k+1)/2,(k+2)/2\}\), for either value of \(\delta\).
Duplication therefore yields
\begin{equation}
 \frac{\Gamma(b)\Gamma(1-a)}
 {\Gamma(b+z/2)\Gamma(1-a+z/2)}
 =2^z\frac{\Gamma(k+1)}{\Gamma(k+1+z)}.
 \label{eq:bridge-gamma-duplication}
\end{equation}
Substituting \eqref{eq:bridge-gamma-half-integer} or
\eqref{eq:bridge-gamma-integer}, followed by
\eqref{eq:bridge-gamma-duplication}, into
\eqref{eq:bridge-normalized-intermediate} proves the theorem.
\end{proof}

\begin{remark}
The statement concerns analytic germs and local logarithms.
It requires no claim about a global zero-free region.
For a complex character the functional equation involves
\(\overline\chi\), and therefore the reflected Taylor coefficients
are conjugated. Replacing \(B^\#\) by \(B\) is justified for a real
character but generally false for a complex one.
\end{remark}

\subsection{Closed logarithmic coefficients at every order}

Choose the normalized local logarithms, each zero at the origin,
of \(L(k+1+z,\chi)/L(k+1,\chi)\) and \(B(t)/B(0)\).
For \(r\geq1\), set
\[
 \lambda_r=
 \left.\frac{d^r}{ds^r}\log L(s,\chi)\right|_{s=k+1},
 \qquad
 b_r=
 \left.\frac{d^r}{dt^r}\log B(t)\right|_{t=0}.
\]
These derivatives do not depend on a choice of logarithm branch.
Write
\[
 H_k^{(r)}=\sum_{j=1}^{k}j^{-r},\qquad H_k=H_k^{(1)},
 \qquad
 \eta(r)=(1-2^{1-r})\zeta(r)\quad(r\geq2).
\]

\begin{theorem}[All-order logarithmic bridges]
\label{thm:bridge-log-jets}
For every \(r\geq1\),
\begin{equation}
 \lambda_r-(-1)^r\overline{b_r}=\kappa_r,
 \label{eq:bridge-log-jets}
\end{equation}
where
\begin{equation}
 \kappa_r=
 \begin{cases}
 \displaystyle\gamma+\log\frac{2\pi}{q}-H_k,
       &r=1,\\[4pt]
 \displaystyle(r-1)!\bigl(\zeta(r)-H_k^{(r)}\bigr),
       &r\geq3\ \text{odd},\\[4pt]
 \displaystyle(r-1)!\bigl(H_k^{(r)}+(-1)^\nu\eta(r)\bigr),
       &r\geq2\ \text{even}.
 \end{cases}
 \label{eq:bridge-kappa}
\end{equation}
\end{theorem}

\begin{proof}
Taking normalized logarithms of
\eqref{eq:bridge-normalized-germ} gives
\[
 \log\frac{L(k+1+z,\chi)}{L(k+1,\chi)}
 =\log K_\nu(z)
  +\overline{\log\frac{B(-\overline z)}{B(0)}}.
\]
The \(r\)-th derivative of the last term at zero is
\((-1)^r\overline{b_r}\). It remains to compute the
derivatives of \(\log K_\nu\).

The convergent gamma expansion, on \(|z|<1\), is
\begin{equation}
 \log\frac{\Gamma(k+1+z)}{\Gamma(k+1)}
 =(H_k-\gamma)z+
 \sum_{r\geq2}\frac{(-1)^r}{r}
       \bigl(\zeta(r)-H_k^{(r)}\bigr)z^r.
 \label{eq:bridge-log-gamma-series}
\end{equation}
The sine and cosine products give, in the same disk,
\begin{align}
 \log\sec\frac{\pi z}{2}
 &=\sum_{m\geq1}
   \frac{(1-2^{-2m})\zeta(2m)}{m}z^{2m},
 \label{eq:bridge-log-sec-series}\\
 \log\frac{\pi z/2}{\sin(\pi z/2)}
 &=\sum_{m\geq1}
   \frac{2^{-2m}\zeta(2m)}{m}z^{2m}.
 \label{eq:bridge-log-sinc-series}
\end{align}
These identities also follow directly by taking logarithms
of the absolutely locally convergent products
\(\cos(\pi z/2)=\prod_{j\geq0}(1-z^2/(2j+1)^2)\) and
\(\sin(\pi z/2)/(\pi z/2)
=\prod_{j\geq1}(1-z^2/(2j)^2)\).

The linear term of \(\log K_\nu\) is the stated
\(\kappa_1\). Odd coefficients of order at least three come
only from the negative of
\eqref{eq:bridge-log-gamma-series}, giving the second case.
For \(r=2m\), the secant coefficient combines with the gamma
coefficient to give
\[
 (r-1)!\left[
  H_k^{(r)}-\zeta(r)+2(1-2^{-r})\zeta(r)\right]
 =(r-1)!\bigl(H_k^{(r)}+\eta(r)\bigr).
\]
Replacing the secant factor by the sine quotient instead gives
\[
 (r-1)!\left[
  H_k^{(r)}-\zeta(r)+2^{1-r}\zeta(r)\right]
 =(r-1)!\bigl(H_k^{(r)}-\eta(r)\bigr).
\]
These are exactly the two even cases in
\eqref{eq:bridge-kappa}.
\end{proof}

At a nonzero reflected value, \(\nu=0\), the first two identities
read
\begin{align}
 \frac{L'(k+1,\chi)}{L(k+1,\chi)}
 +\overline{\frac{L'(-k,\chi)}{L(-k,\chi)}}
 &=\gamma+\log\frac{2\pi}{q}-H_k,
 \label{eq:bridge-first-nonzero}\\
 (\log L)''(k+1,\chi)
 -\overline{(\log L)''(-k,\chi)}
 &=H_k^{(2)}+\frac{\pi^2}{12}.
 \label{eq:bridge-second-nonzero}
\end{align}
At a trivial zero, \(\nu=1\), the Taylor expansion
\[
 B(t)=L'(-k,\chi)
   +\frac{L''(-k,\chi)}2t
   +\frac{L'''(-k,\chi)}6t^2+\cdots
\]
instead gives
\begin{align}
 \frac{L'(k+1,\chi)}{L(k+1,\chi)}
 +\frac12\overline{\frac{L''(-k,\chi)}{L'(-k,\chi)}}
 &=\gamma+\log\frac{2\pi}{q}-H_k,
 \label{eq:bridge-first-zero}\\
 (\log L)''(k+1,\chi)
 -\overline{\left(
    \frac{L'''(-k,\chi)}{3L'(-k,\chi)}
   -\frac{L''(-k,\chi)^2}{4L'(-k,\chi)^2}
   \right)}
 &=H_k^{(2)}-\frac{\pi^2}{12}.
 \label{eq:bridge-second-zero}
\end{align}
These formulas preserve the familiar first-order harmonic
bridge and supply a uniform extension to every derivative
order and both character parities.

\subsection{Raw jets and an exact triangular implementation}

The normalized germ also gives direct formulas for ordinary
derivatives, without first converting each order to logarithmic
derivatives. Define \(K_m\) by
\[
 K_\nu(z)=\sum_{m\geq0}K_m\frac{z^m}{m!},\qquad K_0=1.
\]
Since \(\log K_\nu(z)=\sum_{r\geq1}\kappa_rz^r/r!\),
these coefficients are the complete exponential Bell polynomials:
\begin{equation}
 K_m=
 m!\!\!\sum_{\substack{m_1,\ldots,m_m\geq0\\
                     m_1+2m_2+\cdots+mm_m=m}}
 \prod_{r=1}^{m}\frac1{m_r!}
       \left(\frac{\kappa_r}{r!}\right)^{m_r}
 \qquad(m\geq1).
 \label{eq:bridge-bell-coefficients}
\end{equation}
Leibniz's rule in \eqref{eq:bridge-normalized-germ} proves
\begin{equation}
 \frac{L^{(m)}(k+1,\chi)}{L(k+1,\chi)}
 =
 \sum_{j=0}^{m}\binom mj K_{m-j}(-1)^j
       \overline{\frac{B^{(j)}(0)}{B(0)}}.
 \label{eq:bridge-raw-jets}
\end{equation}
The required regularized derivatives are explicitly
\begin{equation}
 \frac{B^{(j)}(0)}{B(0)}
 =\frac{j!\,\nu!}{(j+\nu)!}
    \frac{L^{(j+\nu)}(-k,\chi)}{L^{(\nu)}(-k,\chi)}.
 \label{eq:bridge-regularized-raw-jets}
\end{equation}
The diagonal coefficient in
\eqref{eq:bridge-raw-jets} is \((-1)^m\), so the transformation
is triangular and invertible. Its coefficients belong to the
ring generated by the explicit \(\kappa_r\)'s. This supplies
exact relation certificates at arbitrary order and avoids
repeated symbolic differentiation across gamma poles.

\begin{remark}[Correction of the complex-character formulas]
The older derivative-relations report omits the conjugation
in versions of the trivial-zero first bridge and the nonzero
second bridge. Equations~\eqref{eq:bridge-first-zero} and
\eqref{eq:bridge-second-nonzero} give the required replacements.
The two main Stieltjes articles already include conjugation
in their displayed first bridges; those correct formulas
should be retained. Exact source locators and numerical
counterexamples to the omitted-conjugation versions appear
in the correction register.
\end{remark}
